%Copyright 2014 Jean-Philippe Eisenbarth
%This program is free software: you can
%redistribute it and/or modify it under the terms of the GNU General Public
%License as published by the Free Software Foundation, either version 3 of the
%License, or (at your option) any later version.
%This program is distributed in the hope that it will be useful,but WITHOUT ANY
%WARRANTY; without even the implied warranty of MERCHANTABILITY or FITNESS FOR A
%PARTICULAR PURPOSE. See the GNU General Public License for more details.
%You should have received a copy of the GNU General Public License along with
%this program.  If not, see <http://www.gnu.org/licenses/>.

%Based on the code of Yiannis Lazarides
%http://tex.stackexchange.com/questions/42602/software-requirements-specification-with-latex
%http://tex.stackexchange.com/users/963/yiannis-lazarides
%Also based on the template of Karl E. Wiegers
%http://www.se.rit.edu/~emad/teaching/slides/srs_template_sep14.pdf
%http://karlwiegers.com
\documentclass{scrreprt}
\usepackage{listings}
\usepackage{underscore}
\usepackage{graphicx}
\usepackage{svg}
\usepackage[bookmarks=true]{hyperref}
\usepackage[utf8]{inputenc}
\usepackage[english]{babel}
\usepackage{hyperref}
\hypersetup{
    bookmarks=false,    % show bookmarks bar?
    pdftitle={Software Requirement Specification},    % title
    pdfauthor={VER VALEM Willian BAYOL Elmer Liliana Lopez Farfan TAGNAOUTI Khaoula},                     % author
    pdfsubject={TeX and LaTeX},                        % subject of the document
    pdfkeywords={TeX, LaTeX, graphics, images}, % list of keywords
    colorlinks=true,       % false: boxed links; true: colored links
    linkcolor=blue,       % color of internal links
    citecolor=black,       % color of links to bibliography
    filecolor=black,        % color of file links
    urlcolor=purple,        % color of external links
    linktoc=page            % only page is linked
}%
\def\myversion{1.0 }
\date{}
%\title{%

%}
\usepackage{hyperref}
\begin{document}

\begin{flushright}
    \rule{16cm}{5pt}\vskip1cm
    \begin{bfseries}
        \Huge{SOFTWARE REQUIREMENTS\\ SPECIFICATION}\\
        \vspace{1.9cm}
        for\\
        \vspace{1.9cm}
        Dominion Logs Datamining\\
        \vspace{1.9cm}
        Prepared by \\ VER VALEM Willian \\ BAYOL Elmer \\ Liliana Lopez Farfan
        \\ TAGNAOUTI Khaoula\\
        \vspace{1.9cm}
        \today\\
    \end{bfseries}
\end{flushright}

\tableofcontents


% \chapter*{Revision History}

% \begin{center}
%     \begin{tabular}{|c|c|c|c|}
%         \hline
%       Name & Date & Reason For Changes & Version\\
%         \hline
%       21 & 22 & 23 & 24\\
%         \hline
%       31 & 32 & 33 & 34\\
%         \hline
%     \end{tabular}
% \end{center}

\chapter{Introduction}

\section{Purpose}
% Identify the product whose software requirements are specified in this
% document, including the revision or release number. Describe the scope of the
% product that is covered by this SRS, particularly if this SRS describes only
% part of the system or a single subsystem.
This Document has as purpose to clarify and describe as precise as possible the
requirements for the development of the program \textit{Dominion DataMining}.\\
\textit{Dominion DataMining} is a program which consists of 3 parties:

\begin{itemize}
  \item \textbf{Parser} \- Responsible to read and remodel the data to be stored.
  \item \textbf{Database Server} \- Store the parsed data
  \item \textbf{Data analyzer} \- Responsible to query the database and generate
    human readable statistics and and plot charts.
\end{itemize}

This requirements Specification is based on the version 0.0.1 of the program.



 \section{Intended Audience and Reading Suggestions}
% Describe the different types of reader that the document is intended for,
% such as developers, project managers, marketing staff, users, testers, and
% documentation writers. Describe what the rest of this SRS contains and how it is
% organized. Suggest a sequence for reading the document, beginning with the
% overview sections and proceeding through the sections that are most pertinent to
% each reader type.
This document is intended for any individual user,developer,tester,project manager or
documentation writer that needs to understand the basic system architecture and its specifications.\\
Here are the potential uses for each one of the reader types:\\
\begin{itemize}

\item \textbf{Developer:}The developer who wants to read,change,modify or add new requirements into
the existing program,must firstly consult this document and update the requirements with
appropriate manner so as to not destroy the actual meaning of them and pass the information
correctly to the next phases of the development process.\\

\item \textbf{User:}The user of this program reviews the diagrams and the specifications presented in this
document and determines if the software has all the suitable requirements and if the software
developer has implemented all of them.\\

\item \textbf{Tester:}The tester needs this document to validate that the initial requirements of this
programs actually corresponds to the executable program correctly.\\

\end{itemize}

\section{Project Scope}
% Provide a short description of the software being specified and its purpose,
% including relevant benefits, objectives, and goals. Relate the software to
% corporate goals or business strategies. If a separate vision and scope document
% is available, refer to it rather than duplicating its contents here.
  The \textit{Dominion} card game has been running on a server from October 2010 to March 2013. The games played on this server have been recorded in \textit{logs} that keep track of every players actions; Those \textit{logs} have been made available to the public. But some information relative to decisions made by the players haven't been recorded.
  \newline A wiki about the game has been created, it offers expert advice in decision making during the game: it mainly offers advice on a strategy level. The goal of this project, submitted by Yvan Le Borgne, researcher at the Labri, is to process the \textit{logs} (more than 12 million games) in order to answer several questions. It mainly consists in comparing the gathered data to the wiki's information in order to validate the usefulness of the ideas and informations given by experts. We will at least create validation tools offering enough efficiency.
  In order to fulfill this objective, the project's goal is to create a database containing all the games played and the to produce an analysis based on this information.
  To do so, our program will first have to extract the \textit{log} data because they are too voluminous and too complicated to use directly (non-standardized structure). We will have to set up data structures allowing us to have a better data storage. A simple way to represent this data will be offered to the user.
  Some information is also missing in the \textit{logs}(for example: cards drawn during a turn). We will try to find a method to rebuild this information and to add it to the existing information.
  We will also try to optimize the opeations made by the data analyzing by finding a balance between memorizing requests and re-computing.
  Finally, we will have to decide wich strategy is used with each game, or even recognize strategy changes during a game. We could also recognize ho applied a strategy first (who \textit{inveted} it).
  If we can detect enough strategy related data, we can also detect the usual learning process of the players.


%\section{References}
% $<$List any other documents or Web addresses to which this SRS refers. These may
% include user interface style guides, contracts, standards, system requirements
% specifications, use case documents, or a vision and scope document. Provide
% enough information so that the reader could access a copy of each reference,
% including title, author, version number, date, and source or location.$>$


\chapter{Overall Description}

\section{Product Perspective}
% $<$Describe the context and origin of the product being specified in this SRS.
% For example, state whether this product is a follow-on member of a product
% family, a replacement for certain existing systems, or a new, self-contained
% product. If the SRS defines a component of a larger system, relate the
% requirements of the larger system to the functionality of this software and
% identify interfaces between the two. A simple diagram that shows the major
% components of the overall system, subsystem interconnections, and external
% interfaces can be helpful.$>$
The \textit{Dominion DataMining} has as objective to be able to read the




\section{Product Functions}
% $<$Summarize the major functions the product must perform or must let the user
% perform. Details will be provided in Section 3, so only a high level summary
% (such as a bullet list) is needed here. Organize the functions to make them
% understandable to any reader of the SRS. A picture of the major groups of
% related requirements and how they relate, such as a top level data flow diagram
% or object class diagram, is often effective.$>$
Following is a list of the expected functions to be available on the software.
\begin{itemize}
  \item{\textbf{Data modeling:}}
    \begin{itemize}
      \item{parser}
      \item{log decompression}
      \item{create object game-log}
      \item{restore game-log}
    \end{itemize}
  \item{\textbf{Data Store:}}
    \begin{itemize}
      \item{Create SQL Database}
      \item{Connect to SQL Database}
      \item{Create NO-SQL Database}
      \item{Connect to NO-SQL Database}
      \item{Compression}
      \item{Load game-log}
      \item{Save game-log}
      \item{Query Interface}
    \end{itemize}
  \item{\textbf{Data Analysis:}}
    \begin{itemize}
      \item{Generate Queries}
      \item{Analyze}
      \item{Plotting}
      \item{Calculate \textit{\textbf{ELO}}}
      \item{Recognize strategies}
      \item{Recognize Greening}
    \end{itemize}
\end{itemize}
\textit{ A more precise description of each functionality will be provided in Section 4 of this document. }


\section{User Classes and Characteristics}
% $<$Identify the various user classes that you anticipate will use this product.
% User classes may be differentiated based on frequency of use, subset of product
% functions used, technical expertise, security or privilege levels, educational
% level, or experience. Describe the pertinent characteristics of each user class.
% Certain requirements may pertain only to certain user classes. Distinguish the
% most important user classes for this product from those who are less important
% to satisfy.$>$

\textbf{Physical Actors:}\\
\begin{itemize}
 \item \textbf{User:} The user will be the one that demands analyses to be done
   on the stored data.
\end{itemize}
\textbf{System Actors:}\\
\begin{itemize}
 \item \textbf{Parser:} The parser is the system has write access to the data
   base and is responsible to store the parsed data on the database.
 \item \textbf{Analyzer:} The Analyzer is the system that will query the data
   base to collect data for analyze , and be able to create persistent queries.
\end{itemize}

\textit{The Software to be developed is aimed to a personal use, so there is no privileged access, the user will have access to the interety of the code and data.}
\section{Operating Environment}
% $<$Describe the environment in which the software will operate, including the
% hardware platform, operating system and versions, and any other software
% components or applications with which it must peacefully coexist.$>$
The Software will be developed for the operational system \textit{Linux} and no
other platform is necessary.
And the software has to run on hardware found on simple personal computers.

\section{Design and Implementation Constraints}
% $<$Describe any items or issues that will limit the options available to the
% developers. These might include: corporate or regulatory policies; hardware
% limitations (timing requirements, memory requirements); interfaces to other
% applications; specific technologies, tools, and databases to be used; parallel
% operations; language requirements; communications protocols; security
% considerations; design conventions or programming standards (for example, if the
% customer’s organization will be responsible for maintaining the delivered
% software).$>$
This section presents an overview of possible issues that cam limit the
development of the Software:\\
The data to be analyzed is considerable big (400Gb) and the client will not provide a
server to work with the data uncompressed what can impact on the performance of
the software.\\
The \textit{Logs} provided by the client don't have an standard format and that
can delay the time to develop the parser, as the problems that it can raise
during the development ,is unknown until the moment.




\section{User Documentation}
% $<$List the user documentation components (such as user manuals, on-line help,
% and tutorials) that will be delivered along with the software. Identify any
% known user documentation delivery formats or standards.$>$
The software will be delivered with a documentation of the source code, the
source code and a user manual.



\chapter{The Logs}

\subsection{Data description}
The data that was provided by the client is compressed in format tar.bz2 , one file for each day of \textit{Log}, the total size of the compressed data is
13Gb.\\
The decompressed data makes a total of 400Gb.\\
Each \textit{Log} file is in \textit{HTML} format.\\
A \textit{Log} should contain:
\begin{itemize}
  \item{\textbf{header}} containing the game number and a winner.
  \item{\textbf{match resume}} containing the cards used on the match and how the match
    finished
  \item{\textbf{player resume}} containing each player victory cards, deck and points.UseCaseDataMining
  \item{\textbf{game log}} that part contains all detailed moves made by the players
    during the match
\end{itemize}

\subsubsection{Data inconsistency}
The \textit{Logs} present some data inconsistency that can be problematic on the
development of the parser.
some problems found at the moment that this SRS is being write are:
\begin{itemize}
  \item \textbf{Log number}\\
    The numeration of the logs aren't unique what implies that can't be used to
    identify a log.
  \item \textbf{User Names}\\
    The logs show that no restriction where made on the user names of the game
    and some user names have Keywords and special characters that can conflict
    with the parser.
   \item \textbf{Missing data}\\
     Some of the logs are missing part of the data (like: the header, users
     resume $\ldots$).
   \item \textbf{Log format}\\
     The syntax used to write the logs are not consistent and present some
     difference between logs.
    \item \textbf{Data compression}\\
      As the data is compressed and the equipment provided to work on the
      project can not handle the data uncompressed the project will have to work
      with the compressed data, and after a first benchmark, 4 hours was needed
      to decompress.
\end{itemize}

\chapter{Functional Requirements}
% $<$This template illustrates organizing the functional requirements for the
% product by system features, the major services provided by the product. You may
% prefer to organize this section by use case, mode of operation, user class,
% object class, functional hierarchy, or combinations of these, whatever makes the
% most logical sense for your product.$>$
The functional requirements will be
% \textbf{the following image illustrate an overview of the use case of the software.}\\
% \includegraphics[width=\textwidth,height=\textheight,keepaspectratio]{UseCaseParsing}
% \newpage
% \textbf{this image shows with more details the use-case of the data analyzes and
% data recovery.}\\
% \\
% \includegraphics[width=\textwidth,height=\textheight,keepaspectratio]{UseCaseDataMining}

% \newpage
% \includegraphics[width=\textwidth,height=\textheight,keepaspectratio]{UseCaseParser}




\section{Data Modeling}

\subsection{Parser::REQ1}

\subsubsection{Description and Priority}

\textbf{Priority level=high}\\

Provided the issues revealed by the logs (like: the inconsistency of the syntax)
a use of a parser made with tools like \textit{YACC and LEX} is not recommended.
For that a parser will be done by using the already present \textit{HTML} and
key words presents on the logs.\\
The parser is responsible for reading the logs and collect the important
information without loosing the structure of the log.\\
A glimpse of the data to be recognized by the parser is illustrated in this
diagram (some data to be determined as all the possibilities inside a log
haven't being revealed yet):\\
\includegraphics[scale=0.35,keepaspectratio]{UseCaseParser}

\subsubsection{Stimulus/response}

This feature will be started by the user via command: parse.
\subsubsection{Dependencies}

depends on: REQ2, REQ3.





\subsection{Log decompression::REQ2}
\subsubsection{Description and Priority}
\textbf{Priority level=high}\\

To be able to work with the large volume of the data a system to decompress the
the logs on demand will be implemented.

This feature is responsible to decompress the tar.bz2 into a temporary folder
and clean the folder.

\subsubsection{Stimulus/response}
This feature will be called by the system feature parser.


\subsection{Create game-log::REQ3}
\subsubsection{Description and Priority}
\textbf{Priority level=high}\\
This functionality is the creation of a data structure where the log data will
be remodeled into.
here is a graphical representation of an overview of the data structure:
\includegraphics[scale=0.5,keepaspectratio]{game-log}


\subsection{restore game-log::REQ4}
\subsubsection{Description and Priority}
\textbf{Priority level=medium}\\
This feature is responsible to restore the missing information on the parsed
logs, by using deduction based on the data obtained from the log.
\subsubsection{Stimulus/response}
This system functionality will be started automatically after a log have being parsed.
\subsubsection{Dependencies}
depends on: REQ1, REQ2 , REQ3.

\newpage
\section{Data Store}

\subsection{Create SQL Database::REQ5}

\subsubsection{Description and Priority}
\textbf{Priority level=high}\\
This functionality is responsible for the creation of the SQL database and the
tables.
for a better understand here is a  sketch of the database structure:\\

\includegraphics[keepaspectratio]{SQL}


\subsection{Connect to SQL Database::REQ6}
\subsubsection{Description and Priority}
\textbf{Priority level=high}\\
This module will be in charge of managing the connection of the application with
the SQL database
\subsubsection{Dependencies}
depends on: REQ5



\subsection{Create NO-SQL Database::REQ7}

\subsubsection{Description and Priority}
\textbf{Priority level=high}\\
This module will be responsible to create the NO-SQL (\textit{Documented based}) database.
as a document based database will be used no data structure has to be defined.

\subsection{Connect to NO-SQL Database::REQ8}
\subsubsection{Description and Priority}
\textbf{Priority level=high}\\
This module will be in charge of managing the connection of the application with
the NO-SQL database
\subsubsection{Dependencies}
depends on: REQ7


\subsection{compression::REQ9}
\subsubsection{Description and Priority}
\textbf{Priority level=medium}\\
The NO-SQL data base will hold most of the data and some level of compression
will have to be applied.
The database in focus to be used at the moment is \textit{mongodb} and it offers
2 levels of compression \textit{Zlib} and \textit{Snappy}.
During the development the \textit{Snappy} will be used as it offers a better
performance.
But the final program should give the user the ability to chose the compression
level.
\subsubsection{Stimulus/response}
The compression will be set on at the time of the creation of the data base.

\subsubsection{Dependencies}
depends on:REQ7

\subsection{Load game-log::REQ10}
\subsubsection{Description and Priority}
\textbf{Priority level=high}\\
This is the functionality that will charge the game-log from the database into
the data structure (see REQ3).
\subsubsection{Dependencies}
depends on: REQ3,REQ7,REQ5.


\subsection{Save game-log::REQ11}
\subsubsection{Description and Priority}
\textbf{Priority level=high}\\
This is the functionality that will transform the game-log int a document
compatible with the database and save it.
The document format used by the database is json.
For more information visit: \url{https://docs.mongodb.org/manual/core/document/}
\subsubsection{Dependencies}
depends on: REQ3,REQ7,REQ5.

\subsection{Query interface::REQ12}
\subsubsection{Description and Priority}
\textbf{Priority level=high}\\
As 2 different kind of databases will be used to store the data.
a query interface that takes the 2 databases in consideration will have to be
implemented.
The query interface will have to choose automatically on what databases the
information will be queried and return the results.
\subsubsection{Dependencies}
depends on: REQ3,REQ7,REQ5.

\newpage



\section{Data Analysis}

\subsection{Generate Queries::REQ13}
\subsubsection{Description and Priority}
\textbf{Priority level=high}\\
This feature will have to be able to generate queries based on a set of
parameters determined by the user.
\subsubsection{Dependencies}
depends on: REQ12,REQ3,REQ7,REQ5.


\subsection{Analyze::REQ14}
\subsubsection{Description and Priority}
\textbf{Priority level=high}\\
This module will get the queries results analyze the data and process the
information generating related human readable statistics data.
\subsubsection{Dependencies}
depends on: REQ13,REQ12,REQ3,REQ7,REQ5.

\subsection{Plotting::REQ15}
\subsubsection{Description and Priority}
\textbf{Priority level=high}\\
This module will receive the statistics generated by REQ14 and generated
graphical charts.
\subsubsection{Stimulus/response}
This feature will be launched by the user via command: graphgem.
\subsubsection{Dependencies}
depends on: REQ14,REQ13,REQ12,REQ3,REQ7,REQ5.



\subsection{Calculate \textit{\textbf{ELO}}::REQ16}
\subsubsection{Description and Priority}
\textbf{Priority level=high}\\
The analyzer will have to calculate the final ELO for each player and what was a
player ELO on a given game.
More information about ELO cam be found on \url{https://en.wikipedia.org/wiki/Elo_rating_system}
\subsubsection{Dependencies}
depends on: REQ12,REQ3,REQ7,REQ5.

\subsection{Recognize strategies::REQ17}
\subsubsection{Description and Priority}
\textbf{Priority level=medium}\\

The dominion wiki describe some strategies that can be used in the game .\\ For exemple :\\
Big Money\\
Beyond Silver \\
Penultimate Province Rule \\
More on \url{http://wiki.dominionstrategy.com/index.php/Strategy}\\

And the analyzer will have to be able to recognize what strategy was used on a
given match (if any strategy was used) and generate statistics.\\
\newpage
\textbf{This image describes the decision process involved in the Big Money strategy}\\
\includegraphics[width=\textwidth,height=\textheight,keepaspectratio]{big-money}\\
The analyzer has to recognize when a player buys only money cards and specific cards related to the big money strategy. It also has to keep track of the remaining provinces to be bought.\\
More on \url{http://wiki.dominionstrategy.com/index.php/Big_Money}\\

In order to recognize the strategy coined Beyond Silver, the analyzer has to recognize when certain type of cards are bought, the list of these cards can be found on \url{http://wiki.dominionstrategy.com/index.php/Silver#Beyond_Silver}\\

In order to recognize that the Penultimate Province Rule is respected (as explained on \url{http://dominionstrategy.com/2011/03/28/the-penultimate-province-rule/}), the analyzer has to keep track of the amount of Victory Points of each player at every turn.

\subsubsection{Dependencies}
depends on: REQ12,REQ3,REQ7,REQ5.

\subsection{Recognize Greening::REQ18}
\subsubsection{Description and Priority}
\textbf{Priority level=high}\\
The analyzer has to be capable of recognize the \textit{\textbf{greening}} moment on each match.
More about \textit{\textbf{greening}} on \url{http://wiki.dominionstrategy.com/index.php/Greening}.
\subsubsection{Dependencies}
depends on: REQ12,REQ3,REQ7,REQ5.


\chapter{Nonfunctional Requirements}

\section{Performance Requirements}
% $<$If there are performance requirements for the product under various
% circumstances, state them here and explain their rationale, to help the
% developers understand the intent and make suitable design choices. Specify the
% timing relationships for real time systems. Make such requirements as specific
% as possible. You may need to state performance requirements for individual
% functional requirements or features.$>$
No specific performance requirement was made by the client. But for the analyzer
a time of running should be less than an hour.



\section{Reliability}
The client demands that the maximum number of logs should be parsed , as the
data provided has some inconsistencies and some logs wont be possible to parse.
the client agrees that between 5 to 10\% of the logs can be ignored.\\

The parser will have to parse 100\% of the data present on the log and Restore
any missing information as possible.\\

\section{Software Quality Attributes}
% $<$Specify any additional quality characteristics for the product that will be
% important to either the customers or the developers. Some to consider are:
% adaptability, availability, correctness, flexibility, interoperability,
% maintainability, portability, reliability, reusability, robustness, testability,
% and usability. Write these to be specific, quantitative, and verifiable when
% possible. At the least, clarify the relative preferences for various attributes,
% such as ease of use over ease of learning.$>$
Even if just a command line is expected for the final program it should have a
easy learn curve and be easy to use.\\
The data generated by the parser don't need to be human readable.\\


%\chapter{Other Requirements}
% $<$Define any other requirements not covered elsewhere in the SRS. This might
% include database requirements, internationalization requirements, legal
% requirements, reuse objectives for the project, and so on. Add any new sections
% that are pertinent to the project.$>$

%\section{Appendix A: Glossary}
% %see https://en.wikibooks.org/wiki/LaTeX/Glossary
% $<$Define all the terms necessary to properly interpret the SRS, including
% acronyms and abbreviations. You may wish to build a separate glossary that spans
% multiple projects or the entire organization, and just include terms specific to
% a single project in each SRS.$>$

%\section{Appendix B: Analysis Models}
% $<$Optionally, include any pertinent analysis models, such as data flow
% diagrams, class diagrams, state-transition diagrams, or entity-relationship
% diagrams.$>$

%\section{Appendix C: To Be Determined List}
% $<$Collect a numbered list of the TBD (to be determined) references that remain
% in the SRS so they can be tracked to closure.$>$

\end{document}
